% ============================================================
% Informe final — Analizador sintáctico LL(1) para un DSL
% Proyecto: Gallinero inteligente — INFO1148, Teoría de la Computación
% Compilar: pdflatex main.tex (dos pasadas para la TOC)
% ============================================================
\documentclass[12pt,a4paper]{article}

% --- Codificación e idioma ---
\usepackage[T1]{fontenc}
\usepackage[utf8]{inputenc}
\usepackage[spanish,es-tabla]{babel}

% --- Márgenes y utilidades ---
\usepackage[margin=2.5cm]{geometry}
\usepackage{booktabs}
\usepackage{longtable}
\usepackage{array}
\usepackage{hyperref}
\usepackage{listings}
\usepackage{xcolor}

\hypersetup{
  colorlinks=true,
  linkcolor=blue!60!black,
  urlcolor=blue!60!black,
  pdftitle={Informe final: Analizador sintáctico LL(1) para un DSL},
  pdfauthor={Tomás Curihual, Rodrigo Gutierrez, Carlos Huenuman},
}

% --- Configuración de bloques de código ---
\definecolor{fondocodigo}{gray}{0.96}
\lstset{
  basicstyle=\ttfamily\small,
  backgroundcolor=\color{fondocodigo},
  breaklines=true,
  columns=fullflexible,
  keepspaces=true,
  frame=single,
  rulecolor=\color{gray!60},
  literate=
    {└──}{{`--}}1
    {├──}{{|--}}1
    {└}{{`}}1
    {├}{{|}}1
    {─}{{-}}1
    {á}{{\'a}}1 {é}{{\'e}}1 {í}{{\'i}}1 {ó}{{\'o}}1 {ú}{{\'u}}1
    {ñ}{{\~n}}1 {Á}{{\'A}}1 {É}{{\'E}}1 {Í}{{\'I}}1 {Ó}{{\'O}}1
    {Ú}{{\'U}}1 {Ñ}{{\~N}}1
    {Σ}{{\ensuremath{\Sigma}}}1
    {ε}{{\ensuremath{$\varepsilon$}}}1,
}

% --- Entornos rápidos ---

\title{ \\
  \vspace{1.5cm}
  {\Huge \textbf{Analizador sintáctico LL(1) para un DSL}} \\
  \vspace{0.4cm}
  {\LARGE Autómata de un gallinero inteligente} \\
  \vspace{2cm}
  {\large Proyecto grupal --- Teoría de la Computación (INFO1148)} \\
  {\large Universidad Católica de Temuco} \\
  \vspace{1.5cm}
}
\author{
  Tomás Curihual \and
  Rodrigo Gutierrez \and
  Carlos Huenuman
}
\date{Octubre de 2026}

\begin{document}

\maketitle

\begin{abstract}
\noindent
Este documento describe el diseño e implementación de un analizador
sintáctico predictivo LL(1) para un lenguaje de dominio específico (DSL)
declarativo orientado a automatizar un gallinero inteligente. El informe
recorre el dominio elegido, la especificación del lenguaje, la gramática
libre de contexto, su análisis (recursión, factorización y ambigüedad), los
conjuntos FIRST y FOLLOW, la tabla predictiva, el parser por descenso
recursivo implementado a mano, la construcción del árbol de sintaxis
abstracta (AST), el manejo de errores sintácticos y la matriz de 46 casos
de prueba. La gramática final es estrictamente LL(1) y toda la
implementación coincide con la especificación documentada en el
repositorio del proyecto.
\end{abstract}

\tableofcontents
\newpage

% =====================================================
\section{Introducción y dominio}
% =====================================================

\subsection{Contexto y problema}
Un gallinero de pequeña escala exige tareas repetitivas y sensibles al
entorno: abrir la puerta por la mañana y cerrarla al anochecer, vigilar la
temperatura y la humedad, reponer alimento y agua, y detectar movimiento de
posibles depredadores. Estas decisiones pueden expresarse como reglas simples
del tipo \emph{«si ocurre una condición, ejecutar una acción»}, lo que las
hace ideales para un lenguaje declarativo.

\subsection{Propósito}
El proyecto diseña e implementa un lenguaje de dominio específico (DSL) y su
analizador sintáctico predictivo \textbf{LL(1)}. El operador escribe reglas
legibles y un controlador las evaluaría cíclicamente para accionar la puerta,
la iluminación, la ventilación, el comedero, el bebedero y la alarma. El
analizador valida que el programa cumpla la sintaxis, construye el árbol de
sintaxis abstracta (AST) y reporta errores con mensajes contextuales que
indican línea, columna y qué se esperaba.

\subsection{Elección del dominio}
Se compararon tres alternativas: un \textbf{gallinero inteligente}, un
\textbf{acuario o estanque} y un \textbf{sistema de riego}. Se eligió el
gallinero inteligente porque:

\begin{itemize}
  \item Tiene sensores y actuadores fáciles de identificar y explicar.
  \item Permite reglas naturales con valores numéricos y booleanos.
  \item Se adapta a la forma \texttt{WHEN <condición> THEN <accion>;}
  \item Tiene relevancia regional y originalidad sin aumentar la complejidad
        sintáctica.
  \item Facilita ejemplos válidos, inválidos y casos límite comprensibles.
\end{itemize}

\subsection{Alcance}
La versión base del lenguaje permite una o más reglas, comparaciones entre
sensores y valores, combinación con \texttt{AND}, \texttt{OR} y \texttt{NOT},
y una acción por regla. La \textbf{extensión creativa} permite varias acciones
por regla separadas por comas, manteniendo compatibilidad total con la base.
Quedan fuera del lenguaje los paréntesis, los comentarios, las acciones con
parámetros, las variables y el control de flujo. La validez sintáctica no
verifica la existencia real de sensores ni la compatibilidad de tipos: eso es
validación semántica y está fuera del alcance.

% =====================================================
\section{Especificación del lenguaje}
% =====================================================

\subsection{Alfabeto}
El alfabeto del lenguaje es:

\begin{lstlisting}
Σ = { A-Z, a-z, 0-9, _, espacio, tab, salto de linea, retorno de carro, >, <, =, !, ;, , }
\end{lstlisting}

No se incluyen paréntesis, puntos, guiones, comillas, caracteres acentuados ni
el símbolo \texttt{\#}. Los caracteres \texttt{=} y \texttt{!} solo aparecen
como parte de los operadores relacionales compuestos (\texttt{==},
\texttt{!=}). Los identificadores son ASCII: se escribe \texttt{lampara}, no
\texttt{lámpara}.

\subsection{Categorías de tokens}
\begin{table}[h!]
\centering
\begin{tabular}{llll}
\toprule
\textbf{Categoría} & \textbf{Tokens} & \textbf{Patrón} & \textbf{Ejemplo} \\
\midrule
Estructural & \texttt{WHEN}, \texttt{THEN} & reservadas & \texttt{WHEN} \\
Lógico      & \texttt{AND}, \texttt{OR}, \texttt{NOT} & reservadas & \texttt{AND} \\
Booleano    & \texttt{TRUE}, \texttt{FALSE} & reservadas & \texttt{TRUE} \\
Identificador & \texttt{ID} & \texttt{[A-Za-z\_][A-Za-z0-9\_]*} & \texttt{humedad} \\
Número      & \texttt{NUM} & \texttt{[0-9]+} & \texttt{100} \\
Relacional  & \texttt{OP\_REL} & \texttt{>= <= == != > <} & \texttt{>=} \\
Separador   & \texttt{COMA} & \texttt{,} & \texttt{,} \\
Terminador  & \texttt{PUNTO\_COMA} & \texttt{;} & \texttt{;} \\
Fin de entrada & \texttt{EOF} & interno & --- \\
Ignorado    & \texttt{ESPACIO} & blancos & --- \\
Error       & \texttt{ERROR} & otros & \texttt{@} \\
\bottomrule
\end{tabular}
\caption{Categorías de tokens del lenguaje.}
\label{tab:tokens}
\end{table}

\subsection{Convenciones léxicas}
\begin{itemize}
  \item Las palabras reservadas se reconocen sin distinguir mayúsculas de
        minúsculas, aunque por convención se escriben en mayúsculas.
  \item Los identificadores sí distinguen mayúsculas y no pueden coincidir
        con una palabra reservada.
  \item Los operadores relacionales se reconocen por la coincidencia más
        larga (\texttt{maximal munch}): \texttt{>=} antes que \texttt{>}.
  \item Un \texttt{=} aislado no es un token válido; la igualdad es
        \texttt{==}. Cualquier símbolo no reconocido produce un token de
        error léxico.
\end{itemize}

\subsection{Palabras reservadas e identificadores del dominio}
El lenguaje tiene siete palabras reservadas: \texttt{WHEN}, \texttt{THEN},
\texttt{AND}, \texttt{OR}, \texttt{NOT}, \texttt{TRUE} y \texttt{FALSE}.

Los identificadores cubren el vocabulario del dominio: sensores numéricos
(\texttt{luz}, \texttt{temperatura}, \texttt{humedad}, \texttt{nivel\_alimento},
\texttt{nivel\_agua}), sensores booleanos (\texttt{movimiento},
\texttt{puerta\_cerrada}, \texttt{ventilacion}, \texttt{alarma\_activa}),
actuadores (\texttt{puerta}, \texttt{lampara}, \texttt{ventilador},
\texttt{calefactor}, \texttt{comedero}, \texttt{bebedero}, \texttt{alarma}),
verbos de acción (\texttt{abrir}, \texttt{cerrar}, \texttt{encender},
\texttt{apagar}, \texttt{activar}, \texttt{desactivar}, \texttt{notificar}) y
el destinatario \texttt{dueno}. Los verbos y destinatarios son
\emph{identificadores}, no palabras reservadas: sintácticamente una acción
solo requiere la forma \texttt{ID ID}.

\subsection{Valores, operadores y separadores}
\begin{itemize}
  \item \texttt{NUM}: entero no negativo de uno o más dígitos. No hay
        negativos, decimales ni cadenas.
  \item Nueve operadores: seis relacionales (\texttt{> < >= <= == !=}) y tres
        lógicos (\texttt{AND}, \texttt{OR}, \texttt{NOT}).
  \item El punto y coma \texttt{;} termina cada regla y es obligatorio.
        La coma \texttt{,} separa acciones dentro de una regla y no puede
        aparecer en ninguna otra posición.
\end{itemize}

\subsection{Precedencia y restricciones de las comparaciones}
\begin{itemize}
  \item \texttt{NOT} tiene mayor precedencia que \texttt{AND}, y
        \texttt{AND} mayor que \texttt{OR}.
  \item Cada comparación tiene exactamente un operador relacional, la forma
        \texttt{ID OP\_REL valor}. No hay comparaciones encadenadas
        (\texttt{luz < 100 < 200}) ni entre dos identificadores.
\end{itemize}

\subsection{Composición de programas}
Un programa contiene una o más reglas; cada regla termina obligatoriamente en
\texttt{;} y tiene al menos una acción. Una acción es un verbo seguido de un
dispositivo, y las reglas con varias acciones las ejecutan de izquierda a
derecha. No se admiten reglas vacías, comas colgantes ni comas consecutivas.
Ejemplo válido:

\begin{lstlisting}
WHEN temperatura > 30 THEN encender ventilador, abrir puerta;
\end{lstlisting}

\subsection{Extensión creativa: múltiples acciones por regla}
La versión base restringía cada regla a una sola acción. La extensión elegida
permite una \emph{rutina} de acciones separadas por comas:

\begin{lstlisting}
WHEN condicion THEN accion, accion, ... ;
\end{lstlisting}

Se evaluaron también los paréntesis, \texttt{ELSE} y las acciones con
duración; se optó por múltiples acciones por su encaje con el dominio, su
bajo impacto en la gramática LL(1) y la facilidad de integrarla en parser,
AST y pruebas. La extensión es retrocompatible: toda regla de la base sigue
siendo válida porque una lista puede tener un solo elemento. Su verificación
LL(1) se detalla en la sección~\ref{sec:tabla-ll1}.

% =====================================================
\section{Gramática libre de contexto}
% =====================================================

\subsection{Formalización}
La gramática se define como la cuádrupla
$G = (N, \Sigma, P, S)$, con $S = $ \texttt{programa} como símbolo inicial.
Los terminales son los tokens del lexer (tabla~\ref{tab:tokens}, incluido
\texttt{EOF}) y los no terminales son:

\begin{lstlisting}
programa, lista_reglas, lista_reglas_prima, regla, condicion, cond_or,
cond_or_prima, cond_and, cond_and_prima, factor, comparacion, valor,
acciones, acciones_prima, accion
\end{lstlisting}

El sufijo \texttt{\_prima} caracteriza la continuación recursiva de cada
regla, técnica habitual para eliminar la recursión por la izquierda.

\subsection{Producciones (BNF)}
\begin{lstlisting}
programa           ::= lista_reglas EOF

lista_reglas       ::= regla lista_reglas_prima

lista_reglas_prima ::= regla lista_reglas_prima
                    | epsilon

regla              ::= WHEN condicion THEN acciones PUNTO_COMA

condicion          ::= cond_or

cond_or            ::= cond_and cond_or_prima

cond_or_prima      ::= OR cond_and cond_or_prima
                    | epsilon

cond_and           ::= factor cond_and_prima

cond_and_prima     ::= AND factor cond_and_prima
                    | epsilon

factor             ::= NOT comparacion
                    | comparacion

comparacion        ::= ID OP_REL valor

valor              ::= NUM
                    | TRUE
                    | FALSE

acciones           ::= accion acciones_prima

acciones_prima     ::= COMA accion acciones_prima
                    | epsilon

accion             ::= ID ID
\end{lstlisting}

La gramática deliberadamente \emph{no} usa las notaciones EBNF
(\texttt{+}, \texttt{*}, \texttt{?}): así puede traducirse directamente al
parser predictivo. La forma compacta equivalente
(\texttt{lista\_reglas\_prima} como \texttt{regla*}, etc.) es equivalente y
se conserva como documentación.

\subsection{Precedencia forzada estructuralmente}
La jerarquía de no terminales fija la precedencia: \texttt{NOT} se aplica a
la comparación que sigue (nivel \texttt{factor}), \texttt{AND} agrupa
factores (nivel \texttt{cond\_and}) y \texttt{OR} es el operador de menor
precedencia (nivel \texttt{cond\_or}).

\subsection{Propiedades}
\begin{itemize}
  \item \textbf{Sin reglas vacías:} \texttt{programa} exige al menos una
        regla, y \texttt{regla} exige al menos una acción.
  \item \textbf{Sin comas colgantes ni consecutivas:} después de
        \texttt{COMA} siempre debe venir una \texttt{accion}.
  \item \textbf{Sin comparaciones encadenadas:} \texttt{comparacion} tiene
        un único \texttt{OP\_REL}.
  \item \textbf{Recursión por la derecha:} toda la repetición usa
        producciones \texttt{\_prima} con recursión derecha, compatible con
        el análisis predictivo LL(1).
  \item \textbf{Derivaciones de ejemplo:} en el repositorio
        (\texttt{docs/gramatica.md}) se documentan las derivaciones completas
        de \texttt{WHEN luz < 100 THEN abrir puerta;}, de una regla con dos
        acciones y de una condición con \texttt{AND}, \texttt{OR} y
        \texttt{NOT}.
\end{itemize}

% =====================================================
\section{Análisis de la gramática}
% =====================================================

\subsection{Recursión por la izquierda}
La gramática original no presenta recursión por la izquierda, ni directa ni
indirecta. Sus producciones recursivas (\texttt{cond\_or\_prima},
\texttt{cond\_and\_prima}, \texttt{acciones\_prima}) ya eran recursivas por
la derecha, y la factorización agregó \texttt{lista\_reglas\_prima} del mismo
tipo. Todo el conjunto es compatible con el análisis predictivo LL(1).

\subsection{Factorización por la izquierda}
La producción \texttt{lista\_reglas} de la gramática previa presentaba un
conflicto FIRST/FIRST:

\begin{lstlisting}
lista_reglas ::= regla lista_reglas
              | regla
\end{lstlisting}

Ambas alternativas comienzan con el no terminal \texttt{regla}, cuyo FIRST es
$\{$\texttt{WHEN}$\}$. Se extrajo el factor común y se introdujo
\texttt{lista\_reglas\_prima}:

\begin{lstlisting}
lista_reglas       ::= regla lista_reglas_prima
lista_reglas_prima ::= regla lista_reglas_prima | epsilon
\end{lstlisting}

Este es el único cambio requerido por la extensión de múltiples acciones
respecto del análisis de la versión base; el resto de la gramática ya estaba
factorizada.

\subsection{Ambigüedad}
La gramática final \textbf{no es ambigua}. La precedencia de los operadores
lógicos está forzada estructuralmente mediante niveles de derivación
(\texttt{NOT} en \texttt{factor}, \texttt{AND} en \texttt{cond\_and},
\texttt{OR} en \texttt{cond\_or}), y la asociatividad queda garantizada por
la recursión por la derecha. Para \texttt{AND} y \texttt{OR}, operadores
matemáticamente asociativos, la forma de agrupación del árbol no altera el
valor de verdad de la condición.

\subsection{Verificación LL(1) de la extensión}
Para la lista de acciones, \texttt{FIRST(acciones\_prima)} =
$\{$\texttt{COMA}, $\epsilon\}$ y \texttt{FOLLOW(acciones\_prima)} =
$\{$\texttt{PUNTO\_COMA}$\}$. Como \texttt{COMA} no pertenece al FOLLOW y
\texttt{PUNTO\_COMA} no pertenece al FIRST, las alternativas tienen FIRST
disjuntos y \textbf{no hay conflicto LL(1)}; tampoco se introduce recursión
izquierda ni ambigüedad nueva.

% =====================================================
\section{FIRST y FOLLOW}
% =====================================================

Los conjuntos se calcularon sobre la gramática factorizada, con recursión por
la derecha ($\epsilon$ indica la producción vacía).

\subsection{Conjuntos FIRST}
\begin{table}[h!]
\centering
\begin{tabular}{ll}
\toprule
\textbf{No terminal} & \textbf{FIRST} \\
\midrule
\texttt{programa}            & \{$\;\texttt{WHEN}\;\}$ \\
\texttt{lista\_reglas}       & \{$\;\texttt{WHEN}\;\}$ \\
\texttt{lista\_reglas\_prima}& \{$\;\texttt{WHEN}, \epsilon\;\}$ \\
\texttt{regla}               & \{$\;\texttt{WHEN}\;\}$ \\
\texttt{condicion}           & \{$\;\texttt{NOT}, \texttt{ID}\;\}$ \\
\texttt{cond\_or}            & \{$\;\texttt{NOT}, \texttt{ID}\;\}$ \\
\texttt{cond\_or\_prima}     & \{$\;\texttt{OR}, \epsilon\;\}$ \\
\texttt{cond\_and}           & \{$\;\texttt{NOT}, \texttt{ID}\;\}$ \\
\texttt{cond\_and\_prima}    & \{$\;\texttt{AND}, \epsilon\;\}$ \\
\texttt{factor}              & \{$\;\texttt{NOT}, \texttt{ID}\;\}$ \\
\texttt{comparacion}         & \{$\;\texttt{ID}\;\}$ \\
\texttt{valor}               & \{$\;\texttt{NUM}, \texttt{TRUE}, \texttt{FALSE}\;\}$ \\
\texttt{acciones}            & \{$\;\texttt{ID}\;\}$ \\
\texttt{acciones\_prima}     & \{$\;\texttt{COMA}, \epsilon\;\}$ \\
\texttt{accion}              & \{$\;\texttt{ID}\;\}$ \\
\bottomrule
\end{tabular}
\caption{Conjuntos FIRST de la gramática factorizada.}
\label{tab:first}
\end{table}

\subsection{Conjuntos FOLLOW}
\begin{table}[h!]
\centering
\begin{tabular}{ll}
\toprule
\textbf{No terminal} & \textbf{FOLLOW} \\
\midrule
\texttt{programa}            & \{$\;\texttt{EOF}\;\}$ \\
\texttt{lista\_reglas}       & \{$\;\texttt{EOF}\;\}$ \\
\texttt{lista\_reglas\_prima}& \{$\;\texttt{EOF}\;\}$ \\
\texttt{regla}               & \{$\;\texttt{WHEN}, \texttt{EOF}\;\}$ \\
\texttt{condicion}           & \{$\;\texttt{THEN}\;\}$ \\
\texttt{cond\_or}            & \{$\;\texttt{THEN}\;\}$ \\
\texttt{cond\_or\_prima}     & \{$\;\texttt{THEN}\;\}$ \\
\texttt{cond\_and}           & \{$\;\texttt{OR}, \texttt{THEN}\;\}$ \\
\texttt{cond\_and\_prima}    & \{$\;\texttt{OR}, \texttt{THEN}\;\}$ \\
\texttt{factor}              & \{$\;\texttt{AND}, \texttt{OR}, \texttt{THEN}\;\}$ \\
\texttt{comparacion}         & \{$\;\texttt{AND}, \texttt{OR}, \texttt{THEN}\;\}$ \\
\texttt{valor}               & \{$\;\texttt{AND}, \texttt{OR}, \texttt{THEN}\;\}$ \\
\texttt{acciones}            & \{$\;\texttt{PUNTO\_COMA}\;\}$ \\
\texttt{acciones\_prima}     & \{$\;\texttt{PUNTO\_COMA}\;\}$ \\
\texttt{accion}              & \{$\;\texttt{COMA}, \texttt{PUNTO\_COMA}\;\}$ \\
\bottomrule
\end{tabular}
\caption{Conjuntos FOLLOW de la gramática factorizada.}
\label{tab:follow}
\end{table}

\subsection{Observaciones}
\begin{itemize}
  \item \texttt{FOLLOW(programa)} $= \{\texttt{EOF}\}$: el análisis termina
        cuando la entrada se agota.
  \item \texttt{FOLLOW(regla)} $= \{\texttt{WHEN}, \texttt{EOF}\}$: después
        de una regla puede comenzar otra regla o terminar el programa.
  \item \texttt{FOLLOW(accion)} $= \{\texttt{COMA},
        \texttt{PUNTO\_COMA}\}$: refleja directamente la extensión de
        múltiples acciones, que decide si la lista continúa o termina.
\end{itemize}

% =====================================================
\section{Tabla LL(1)}
\label{sec:tabla-ll1}
% =====================================================

La tabla predictiva mapea cada par (no terminal, token) a la producción a
aplicar. Se construyó a partir de los conjuntos FIRST y FOLLOW de las
tablas~\ref{tab:first} y~\ref{tab:follow}: una producción se agrega en la
celda $(A, a)$ si $a \in \text{FIRST}(\alpha)$, o en todas las celdas
$(A, b)$ con $b \in \text{FOLLOW}(A)$ si $\alpha \Rightarrow^* \epsilon$.

\begin{table}[h!]
\centering
\begin{tabular}{lll}
\toprule
\textbf{No terminal} & \textbf{Token} & \textbf{Producción} \\
\midrule
\texttt{programa}            & \texttt{WHEN}         & \texttt{programa ::= lista\_reglas EOF} \\
\texttt{lista\_reglas}       & \texttt{WHEN}         & \texttt{lista\_reglas ::= regla lista\_reglas\_prima} \\
\texttt{lista\_reglas\_prima}& \texttt{WHEN}         & \texttt{lista\_reglas\_prima ::= regla lista\_reglas\_prima} \\
\texttt{lista\_reglas\_prima}& \texttt{EOF}          & \texttt{lista\_reglas\_prima ::= epsilon} \\
\texttt{regla}               & \texttt{WHEN}         & \texttt{regla ::= WHEN condicion THEN acciones PUNTO\_COMA} \\
\texttt{condicion}           & \texttt{NOT}, \texttt{ID} & \texttt{condicion ::= cond\_or} \\
\texttt{cond\_or}            & \texttt{NOT}, \texttt{ID} & \texttt{cond\_or ::= cond\_and cond\_or\_prima} \\
\texttt{cond\_or\_prima}     & \texttt{OR}           & \texttt{cond\_or\_prima ::= OR cond\_and cond\_or\_prima} \\
\texttt{cond\_or\_prima}     & \texttt{THEN}         & \texttt{cond\_or\_prima ::= epsilon} \\
\texttt{cond\_and}           & \texttt{NOT}, \texttt{ID} & \texttt{cond\_and ::= factor cond\_and\_prima} \\
\texttt{cond\_and\_prima}    & \texttt{AND}          & \texttt{cond\_and\_prima ::= AND factor cond\_and\_prima} \\
\texttt{cond\_and\_prima}    & \texttt{OR}, \texttt{THEN} & \texttt{cond\_and\_prima ::= epsilon} \\
\texttt{factor}              & \texttt{NOT}          & \texttt{factor ::= NOT comparacion} \\
\texttt{factor}              & \texttt{ID}           & \texttt{factor ::= comparacion} \\
\texttt{comparacion}         & \texttt{ID}           & \texttt{comparacion ::= ID OP\_REL valor} \\
\texttt{valor}               & \texttt{NUM}          & \texttt{valor ::= NUM} \\
\texttt{valor}               & \texttt{TRUE}         & \texttt{valor ::= TRUE} \\
\texttt{valor}               & \texttt{FALSE}        & \texttt{valor ::= FALSE} \\
\texttt{acciones}            & \texttt{ID}           & \texttt{acciones ::= accion acciones\_prima} \\
\texttt{acciones\_prima}     & \texttt{COMA}         & \texttt{acciones\_prima ::= COMA accion acciones\_prima} \\
\texttt{acciones\_prima}     & \texttt{PUNTO\_COMA}  & \texttt{acciones\_prima ::= epsilon} \\
\texttt{accion}              & \texttt{ID}           & \texttt{accion ::= ID ID} \\
\bottomrule
\end{tabular}
\caption{Tabla de análisis predictivo LL(1).}
\label{tab:ll1}
\end{table}

Las celdas vacías (intersecciones no listadas) corresponden a errores
sintácticos. La tabla es \textbf{estrictamente LL(1)}: tras la factorización
no existe ningún par (no terminal, token) con más de una producción, de modo
que el parser siempre puede predecir la derivación correcta con un único
token de anticipación (\emph{lookahead}).

% =====================================================
\section{Diseño e implementación del parser}
% =====================================================

\subsection{Decisiones generales}
El parser se implementa \textbf{a mano}, por \textbf{descenso recursivo},
siguiendo exactamente la tabla de la sección~\ref{sec:tabla-ll1}: cada no
terminal corresponde a un método. No se usa ningún generador automático de
parsers ni bibliotecas externas (el proyecto usa solo la biblioteca estándar
de Python). La implementación reside en \texttt{src/lexer.py} y
\texttt{src/parser.py}. Esta decisión fue deliberada: escribir el descenso
recursivo directamente desde la tabla garantiza la trazabilidad entre
gramática, tabla y código, y obliga a entender cada producción durante la
implementación.

\subsection{Pipeline de análisis}
El proceso tiene tres etapas:
\begin{enumerate}
  \item El \textbf{lexer} (\texttt{src/lexer.py}) convierte la cadena en una
        lista de tokens terminada en \texttt{EOF}. Reconoce reservadas sin
        distinguir mayúsculas, usa \emph{maximal munch} para los operadores
        relacionales, ignora blancos y marca con \texttt{ERROR} los símbolos
        no pertenecientes al alfabeto.
  \item El \textbf{parser} (\texttt{src/parser.py}) recibe la lista de
        tokens y la recorre con un solo token de anticipación
        (\texttt{self.actual}). El método \texttt{consumir(tipo)} verifica
        que el token actual sea del tipo esperado o reporta el error.
  \item \textbf{Concurrente con el análisis}, el parser construye el AST
        (sección~\ref{sec:ast}): los métodos de producción no solo consumen
        tokens, también devuelven nodos.
\end{enumerate}

\subsection{Estructura del módulo parser}
La correspondencia gramática--métodos es directa:
\begin{table}[h!]
\centering
\begin{tabular}{ll}
\toprule
\textbf{Producción} & \textbf{Método en \texttt{Parser}} \\
\midrule
\texttt{programa}            & \texttt{programa()} \\
\texttt{lista\_reglas}       & \texttt{lista\_reglas()} \\
\texttt{lista\_reglas\_prima}& \texttt{lista\_reglas\_prima()} \\
\texttt{regla}               & \texttt{regla()} \\
\texttt{condicion / cond\_or}& \texttt{condicion()} / \texttt{cond\_or()} \\
\texttt{cond\_or\_prima}     & \texttt{cond\_or\_prima(izq)} \\
\texttt{cond\_and}           & \texttt{cond\_and()} \\
\texttt{cond\_and\_prima}    & \texttt{cond\_and\_prima(izq)} \\
\texttt{factor}              & \texttt{factor()} \\
\texttt{comparacion}         & \texttt{comparacion()} \\
\texttt{valor}               & \texttt{valor()} \\
\texttt{acciones}            & \texttt{acciones()} \\
\texttt{acciones\_prima}     & \texttt{acciones\_prima(accs)} \\
\texttt{accion}              & \texttt{accion()} \\
\bottomrule
\end{tabular}
\caption{Correspondencia entre producciones y métodos del parser.}
\label{tab:parser-metodos}
\end{table}

Los métodos \texttt{\_prima} reciben el subárbol construido hasta el momento
y lo extienden hacia la derecha, replicando la recursión derecha de la
gramática y agrupando las operaciones binarias \texttt{AND}/\texttt{OR} de
izquierda a derecha.

\subsection{Uso desde la línea de comandos}
El analizador recibe el programa como argumento:
\begin{lstlisting}
python src/parser.py "WHEN luz < 100 THEN abrir puerta;"
\end{lstlisting}
Devuelve el código \texttt{0} si la entrada es válida (e imprime el árbol),
\texttt{1} ante cualquier entrada inválida (incluidos los símbolos no
reconocidos por el lexer, que llegan al parser como token \texttt{ERROR}), y
\texttt{2} si no se pasa ningún argumento. Estos códigos permiten que las
pruebas por matriz y la suite de \texttt{unittest} verifiquen también el
comportamiento completo de la línea de comandos.

% =====================================================
\section{Árbol sintáctico / AST}
\label{sec:ast}
% =====================================================

\subsection{Diseño de nodos}
El AST se define en \texttt{src/arbol.py} con seis clases: \texttt{NodoAST}
(base), \texttt{Programa} (lista de reglas), \texttt{Regla} (una condición y
una \emph{lista ordenada} de acciones), \texttt{Comparacion}
(\texttt{ID OP\_REL valor}), \texttt{OperacionBinaria} (\texttt{AND},
\texttt{OR}) y \texttt{OperacionUnaria} (\texttt{NOT}). Cada nodo sabe
imprimirse jerárquicamente, lo que permite inspeccionar el resultado del
parser desde la consola.

El diseño es deliberadamente \emph{abstracto}: los nodos del AST no contienen
los delimitadores \texttt{WHEN} y \texttt{THEN} (que son estructura, no
información) y las listas de reglas y acciones se materializan como listas de
Python durante el análisis. Esta información coincide con la gramática: cada
nodo corresponde a un no terminal \emph{significativo} de la gramática de la
sección 3.

\subsection{Ejemplos}
Para la entrada \texttt{WHEN temperatura > 30 THEN encender ventilador, abrir
puerta;} el parser produce (la salida se muestra tal como la imprime el
módulo en consola, con las ramas del árbol):

\begin{lstlisting}
Entrada sintácticamente válida. Árbol generado:
└── Programa
  ├── Regla (WHEN ... THEN ...)
    ├── Condicion:
      └── Comparacion: temperatura > 30
    └── Acciones:
      └── Accion: encender ventilador
      └── Accion: abrir puerta
\end{lstlisting}

Para una condición compuesta como \texttt{WHEN movimiento == TRUE AND NOT
alarma\_activa == TRUE THEN activar alarma;} la salida evidencia la
precedencia estructural descrita en la sección~\ref{sec:tabla-ll1}: el
\texttt{AND} agrupa a nivel superior la comparación y la negación:

\begin{lstlisting}
Entrada sintácticamente válida. Árbol generado:
└── Programa
  ├── Regla (WHEN ... THEN ...)
    ├── Condicion:
      ├── Operador Binario: AND
        └── Comparacion: movimiento == TRUE
        ├── Operador Unario: NOT
          └── Comparacion: alarma_activa == TRUE
    └── Acciones:
      └── Accion: activar alarma
\end{lstlisting}

\subsection{Impresión portable}
La impresión del árbol dibuja las ramas con caracteres de caja Unicode (los
mismos que se observan al ejecutar el programa en una consola compatible) y,
para no depender de la codificación del terminal, el módulo reconfigura sus
flujos de salida con \texttt{errors="replace"}. Así se evita el error
\texttt{UnicodeEncodeError} en terminales Windows configuradas con la página
de códigos \texttt{cp1252} (que no soporta esos símbolos): reproducir los
árboles solo depende de Python, no del terminal del lector.

% =====================================================
\section{Manejo de errores}
% =====================================================

\subsection{Filosofía}
Los errores se centralizan en \texttt{src/errores.py} mediante la excepción
\texttt{ErrorSintactico}. El análisis no continúa tras el primer error
(\emph{panic} simple): se levanta la excepción con la posición exacta del
token problemático y el programa termina con código \texttt{1}. Dado que el
lexer marca con el token \texttt{ERROR} los símbolos ajenos al alfabeto
(declarándolos en la fase léxica) y los entrega al parser, toda la detección
se reporta bajo la interfaz común de \texttt{ErrorSintactico}; de este modo se
conserva una única responsabilidad en mensajes y la suite de pruebas verifica
el comportamiento completo (mensaje, posición y código de salida).

\subsection{Características de los mensajes}
\begin{itemize}
  \item \textbf{Posición precisa:} línea y columna del token que disparó el
        error.
  \item \textbf{Traducción técnica $\to$ humano:} el token se describe en
        vocabulario natural (\texttt{EOF} $\to$ «el final de la entrada»,
        \texttt{ID} $\to$ «el identificador '\emph{lexema}'», etc.).
  \item \textbf{Mensajes contextuales:} el mensaje indica \emph{qué se
        esperaba}: «se esperaba la palabra reservada THEN después de la
        condición», «se esperaba el punto y coma ';' para terminar la regla».
  \item \textbf{Ayudas por símbolo:} los símbolos confusos reciben pistas:
        \texttt{(} $\to$ «los paréntesis no forman parte del lenguaje»,
        \texttt{=} $\to$ «para comparar igualdad use '=='».
  \item \textbf{Sugerencias contextuales del parser:} si después de una
        acción aparece un \texttt{ID}, el parser sugiere «Si hay otra acción,
        sepárela con una coma (',').»; si falta \texttt{THEN} con un
        identificador en curso, sugiere un operador \texttt{AND}/\texttt{OR}.
\end{itemize}

\subsection{Ejemplos}
Los mensajes de la salida de error estándar son literales, obtenidos ejecutando
el analizador tal como se muestra (cada mensaje se imprime en una sola línea;
el salto de línea en el PDF corresponde al ajuste del bloque).

\begin{lstlisting}
> python src/parser.py "WHEN luz < 100 THEN abrir puerta"
Error sintáctico en línea 1, columna 33: se esperaba el punto y coma ';' para terminar la regla; se encontró el final de la entrada.

> python src/parser.py "WHEN luz = 100 THEN abrir puerta;"
Error sintáctico en línea 1, columna 10: se esperaba un operador relacional (>, <, >=, <=, ==, !=); se encontró el símbolo no reconocido '=' (para comparar igualdad use '==').

> python src/parser.py "WHEN luz < 100 abrir puerta;"
Error sintáctico en línea 1, columna 16: se esperaba la palabra reservada THEN después de la condición; se encontró el identificador 'abrir'. Si la condición continúa, falta un operador AND u OR.

> python src/parser.py "WHEN temperatura > 30 THEN encender ventilador abrir puerta;"
Error sintáctico en línea 1, columna 48: se esperaba el punto y coma ';' para terminar la regla; se encontró el identificador 'abrir'. Si hay otra acción, sepárela con una coma (',').

> python src/parser.py ""
Error sintáctico en línea 1, columna 1: la entrada está vacía; se esperaba al menos una regla iniciada por WHEN.

> python src/parser.py
Uso: python src/parser.py "<programa>"
\end{lstlisting}

\subsection{Cobertura}
La suite \texttt{tests/test\_errores.py} contiene \textbf{30 pruebas} que
verifican los mensajes, las posiciones y los códigos de salida ante entradas
inválidas, ejecutándose con \texttt{python tests/test\_errores.py} o como
parte de \texttt{python -m unittest discover -s tests -v}.

% =====================================================
\section{Casos de prueba}
% =====================================================

\subsection{Estructura de la matriz}
La matriz completa tiene \textbf{46 casos}: \textbf{25} entradas válidas y
\textbf{21} inválidas, organizados en 13 casos funcionales válidos (\texttt{V}),
14 funcionales inválidos (\texttt{I}), 9 casos límite (\texttt{L}) y 10 casos
de la extensión creativa (\texttt{E}, 5 válidos y 5 inválidos). Se ejecuta con
\texttt{python tests/test\_parser.py} y se regenera con la opción
\texttt{-{}-tabla tests/casos\_prueba.md}.

En las entradas válidas también se comprueba la \emph{cantidad de reglas y
acciones} del AST, no solo la aceptación. La tabla completa está en
\texttt{tests/casos\_prueba.md} del repositorio; a continuación se muestra
una selección representativa:

\begin{longtable}{llp{7.2cm}}
\caption{Selección representativa de la matriz de 46 casos.}
\label{tab:casos}\\
\toprule
\textbf{ID} & \textbf{Cat.} & \textbf{Entrada / Verificación} \\
\midrule
\endhead
V01 & Válida & \texttt{WHEN luz < 100 THEN abrir puerta;} — regla mínima. \\
V08 & Válida & \texttt{WHEN temperatura > 30 OR ventilacion == FALSE THEN encender ventilador, abrir puerta;} — \texttt{OR} + extensión. \\
V11 & Válida & \texttt{when LUZ < 100 then abrir puerta;} — reservadas en minúsculas. \\
V12 & Válida & \texttt{WHEN temperatura == 000000 THEN encender calefactor;} — número con ceros. \\
V13 & Válida & \texttt{WHEN humedad >= 80 AND ventilacion != TRUE OR movimiento == FALSE THEN ...} — precedencia \texttt{AND} sobre \texttt{OR}. \\
I03 & Inválida & \texttt{WHEN luz < THEN abrir puerta;} — falta el valor. \\
I05 & Inválida & \texttt{WHEN luz < 100 THEN abrir;} — acción sin dispositivo. \\
I08 & Inválida & \texttt{WHEN luz < 100 THEN abrir puerta,, encender lampara;} — comas consecutivas. \\
I10 & Inválida & \texttt{WHEN luz < 100 THEN abrir puerta @;} — símbolo no perteneciente al alfabeto, marcado por el lexer como token \texttt{ERROR}. \\
I12 & Inválida & \texttt{WHEN luz = 100 THEN abrir puerta;} — la igualdad requiere \texttt{==}. \\
L03 & Límite & Cadena profunda de 9 comparaciones con \texttt{AND}. \\
L05 & Límite & Rutina de 5 acciones en una sola regla. \\
L08 & Límite & \texttt{""} — entrada vacía: límite inferior del lenguaje. \\
L09 & Límite & Solo espacios: el lexer produce únicamente \texttt{EOF}. \\
E02 & Extensión & \texttt{WHEN luz < 100 THEN abrir puerta, encender lampara, cerrar ventana;} — rutina de 3 acciones. \\
E09 & Extensión & \texttt{WHEN temperatura > 30 THEN encender ventilador,, abrir puerta;} — dos comas consecutivas. \\
\bottomrule
\end{longtable}

Los casos límite \texttt{L08} y \texttt{L09} representan el límite inferior
del lenguaje: el programa exige al menos una regla, por lo que la entrada
vacía y la entrada con solo espacios deben ser rechazadas.

% =====================================================
\section{Conclusiones}
% =====================================================

\subsection{Resultados principales}
\begin{itemize}
  \item La gramática resultado es \textbf{estrictamente LL(1)}: la factorización
        de \texttt{lista\_reglas} y la recursión derecha en los no terminales
        \texttt{\_prima} eliminaron el único conflicto FIRST/FIRST, y la
        extensión de múltiples acciones se integró sin conflictos porque
        \texttt{COMA} $\notin$ \texttt{FOLLOW(acciones\_prima)}.
  \item La especificación, los documentos de diseño y el código coinciden:
        la tabla predictiva es un reflejo directo de los métodos del parser y
        la matriz de 46 casos se pasa completa, lo que asegura la coherencia
        transversal exigida por la rúbrica.
  \item El enfoque de especificar primero (lenguaje $\to$ gramática $\to$
        análisis $\to$ FIRST/FOLLOW $\to$ tabla) y después implementar hizo
        que el descenso recursivo fuera casi mecánico de escribir y fácil de
        verificar.
  \item La calidad de los mensajes de error agrega valor real al DSL: indicar
        posición exacta, qué se esperaba y ofrecer sugerencias contextuales
        (coma faltante, \texttt{==} en vez de \texttt{=}).
  \item Trabajo futuro: un analizador semántico que valide la existencia de
        sensores y actuadores y la compatibilidad de tipos, la recuperación
        de errores en modo pánico para reportar varios errores por
        ejecución, y un intérprete que ejecute las rutinas de acciones.
\end{itemize}

% =====================================================
\appendix
% =====================================================

\section{Anexos}
% =====================================================

\subsection{Repositorio del proyecto}
\url{https://github.com/HJFlash/Proyecto-TC-2026}

La documentación de diseño asociada, con la que este informe fue redactado y
verificado, se encuentra en el directorio \texttt{docs/} del repositorio:
\begin{itemize}
  \item \texttt{docs/lenguaje.md} — definición del lenguaje y tokens.
  \item \texttt{docs/extension.md} — registro de la extensión creativa.
  \item \texttt{docs/gramatica.md} — gramática libre de contexto y derivaciones.
  \item \texttt{docs/analisis\_gramatica.md} — análisis de la gramática.
  \item \texttt{docs/first\_follow.md} — conjuntos FIRST y FOLLOW.
  \item \texttt{docs/tabla\_ll1.md} — tabla predictiva LL(1).
\end{itemize}

\subsection{Capturas de pantalla}
\textbf{Por completar.} Las capturas de las ejecuciones reales (ejecución
válida con el árbol generado, ejecución inválida con el mensaje de error, y
el paso de las 46 pruebas de la matriz) se tomarán y agregarán en la issue
siguiente. Para copiar las imágenes basta ejecutar los comandos del
anexo~\ref{sec:anexo-reproduccion} y capturar las consolas correspondientes.

\subsection{Cómo reproducir los resultados}
\label{sec:anexo-reproduccion}
\begin{lstlisting}
# Clonar e instalar
git clone https://github.com/HJFlash/Proyecto-TC-2026.git
cd Proyecto-TC-2026
python -m venv .venv && .venv/Scripts/Activate.ps1  # Windows
# (o: source .venv/bin/activate)                     # Linux/macOS

# Matriz completa de 46 casos
python tests/test_parser.py

# Pruebas de manejo de errores (30)
python tests/test_errores.py

# Suite de unittest
python -m unittest discover -s tests -v
\end{lstlisting}

\subsection{Estructura del código}
\begin{lstlisting}
src/
  lexer.py      # Analizador lexico: cadena -> tokens
  parser.py     # Parser LL(1) por descenso recursivo -> AST
  errores.py    # Excepcion y mensajes de error sintactico
  arbol.py      # Nodos del AST e impresion jerarquica
tests/
  casos_prueba.md   # Matriz de 46 casos (generable con --tabla)
  test_parser.py    # Ejecucion de la matriz
  test_errores.py   # 30 pruebas de manejo de errores
\end{lstlisting}

\end{document}
